\documentclass{jsarticle}
\usepackage[dvipdfmx, final]{graphicx}
\usepackage{amsmath}
\usepackage{amssymb}
\begin{document}
\title{}
\author{}
\date{}
\maketitle

$$ {d \over {d f}}(\nabla^2 \psi) = {d \over {d f}}(-{2m \over \hbar^2} V \psi) \\
= -{2m \over \hbar^2} V {d \psi \over d f}$$

大域的積分多様体:
$$\int (\nabla^2 \psi) d x_m = \int (-{2m \over \hbar^2} V \psi) d x_m \\
= -{2m \over \hbar^2} \int V \psi d x_m
$$

AdS5多様体:
大域的微分:
$${d \over {d f}} (e^{-2kz} (\eta_{\mu\nu} d x^{\mu} d x^{\nu} + d z^2)) = e^{-2kz} (\eta_{\mu\nu} {d x^{mu} \over {d f}} d x^{\nu} + {d z^2 \over d f})
$$
大域的積分多様体:
$$
\int e^{-2kz} (\eta_{\mu\nu} d x^{mu} d x^{\nu} + d z^2) d x_m = \int e^{-2kz} (\eta_{\mu\nu} x^{\mu} x^{\nu} + z^2) d x_m
$$

カルーツァ・クライン方程式:
大域的微分:
$${d \over {d f}} (R_{\mu\nu} - {1 \over 2} g_{\mu\nu} R) = {8\pi G \over c^4} {d T_{\mu\nu} \over d f} + {1 \over L^2} {d g_{\mu\nu} \over d f}
$$

大域的積分多様体:
$$\int (R_{\mu\nu} - {1 \over 2} g_{\mu\nu} R) d x_m = {8\pi G \over c^4} \int T_{\mu\nu} d x_m + {1 \over L^2} \int g_{\mu\nu} d x_m
$$

ディラック方程式:
大域的微分:
$$
{d \over {d f}} (i\hbar {\partial \Psi \over \partial t} - mc^2 \Psi) = 0
$$

大域的積分多様体:
$$\int (i\hbar {\partial \Psi \over \partial t} - mc^2 \Psi) d x_m = 0
$$

小林・益川理論:
大域的微分:
$${d \over {d f}} L_{CKM} = {g_W \over \sqrt{2}} {d V_{ij} \over d f} \bar{\Psi_i^L} \gamma^{\mu} W_{\mu} \Psi_j^L + h.c.
$$

大域的積分多様体:
$$\int L_{CKM} d x_m = {g_W \over \sqrt{2}} \int V_{ij} \bar{\Psi_i^L} \gamma^{\mu} W_{\mu} \Psi_j^L d x_m + h.c.
$$
これらの方程式を大域的微分と大域的積分多様体の形式で書き換えることで、より幾何学的な理解が深まると思います。各物理量の微分や積分の形式を明示することで、物理現象と数学構造との対応がより明確になります。
\end{document}
